% from ega4-17-fr.tex
\oldpage[IV-4]{95}

Notons $\alpha$ et $-\beta$ les homomorphismes figurant au second membre de (17.13.5.4), $\gamma$ et $\delta$ ceux figurant au second membre de (17.13.5.2). Dire que la suite (17.13.6.2)
\[
0\to\Omega^1_{X/S}\otimes_{\sh{O}_X}\sh{O}_T
\xrightarrow{\alpha\oplus(-\beta)}
(\Omega^1_{Y/S}\otimes_{\sh{O}_Y}\sh{O}_T)\oplus
(\Omega^1_{Z/S}\otimes_{\sh{O}_Z}\sh{O}_T)
\xrightarrow{\gamma+\delta}\Omega^1_{T/S}\to0
\]
est exacte signifie que, dans la catégorie des $\sh{O}_T$-Modules, $\Omega^1_{X/S}\otimes_{\sh{O}_X}\sh{O}_T$ s'identifie canoniquement au \emph{produit fibré} de $\Omega^1_{Y/S}\otimes_{\sh{O}_Y}\sh{O}_T$ et $\Omega^1_{Z/S}\otimes_{\sh{O}_Z}\sh{O}_T$ sur $\Omega^1_{T/S}$, pour les homomorphismes $\gamma$ et $\delta$. Le même raisonnement que dans (0, 18.1.2 et 18.1.3), où l'on remplace anneaux et idéaux bilatères respectivement par Modules et sous-Modules, fournit un diagramme commutatif
{\scriptsize
\[
\begin{tikzcd}[column sep=small,row sep=large]
0 \ar[dr] & & 0 \ar[d] & 0 \ar[d] & \\
& (\sh{N}_{Y/X}\otimes\sh{O}_T)\oplus(\sh{N}_{Z/X}\otimes\sh{O}_T)
  \ar[r] \ar[d] \ar[dr,"\iota" description]
& \sh{N}_{Z/X}\otimes\sh{O}_T \ar[r,"\sim"] \ar[d]
& \sh{N}_{T/Y} \ar[d] & \\
0 \ar[r] & \sh{N}_{Y/X}\otimes\sh{O}_T \ar[r] \ar[d]
& \Omega^1_{X/S}\otimes\sh{O}_T \ar[r,"\alpha"] \ar[d,"\beta"] \ar[dr,"\varepsilon" description]
& \Omega^1_{Y/S}\otimes\sh{O}_T \ar[r] \ar[d,"\gamma"] & 0\\
0 \ar[r] & \sh{N}_{T/Z} \ar[r]
& \Omega^1_{Z/S}\otimes\sh{O}_T \ar[r,"\delta"] \ar[d]
& \Omega^1_{T/S} \ar[r] \ar[d] \ar[dr] & 0\\
& & 0 & 0 & 0
\end{tikzcd}
\]
}
où les 3\textsuperscript{e} et 4\textsuperscript{e} lignes et colonnes sont exactes en vertu des hypothèses de lissité. Le fait que l'homomorphisme composé $\varepsilon=\gamma\circ\alpha=\delta\circ\beta$ soit l'homomorphisme correspondant à l'injection canonique $T\to X$ résulte de (0, 20.5.2.7). Comme la diagonale du diagramme précédent est exacte, $\operatorname{Ker}(\varepsilon)$ s'identifie canoniquement avec $\sh{N}_{T/X}$, compte tenu de ce que $T$ est lisse sur $S$ (17.2.5); donc on a un isomorphisme canonique (17.13.8.1) réciproque de $\iota$.

\begin{env}[17.13.9]
\label{IV.17.13.9}
On peut généraliser les résultats de (17.13.5) et (17.13.6) à un nombre fini quelconque de $S$-morphismes $f_i:Y_i\to X$ ($i\in I$, $I$ ensemble fini). Pour cela, notons $X^I$ le produit de la famille $(X_i)_{i\in I}$ de $S$-préschémas tous identiques à $X$ (I, 3.3.5), et soient $p_i$ ($i\in I$) les projections canoniques. Le \emph{morphisme diagonal} $\Delta:X\to X^I$ est défini (comme dans (I, 5.3.1)) comme l'unique $S$-morphisme tel que $p_i\circ\Delta=1_X$ pour tout $i\in I$. C'est une $X$-section de $X^I$ pour le morphisme $p_1$, donc (I, 5.3.11) une \emph{immersion}.

